


\subsection{Objetivo}
    YPF ha implementado \emph{cartelería digital en estaciones}, lo que permite efectuar comunicación comercial dinámica. El objetivo de este bloque se centra en evaluar si esta infraestructura impacta en ventas y qué tipo de mensajes son más efectivos. EN lineas generales, se estimara  cómo ciertas características físicas y de infraestructura de las estaciones de servicio, inciden en las ventas de combustibles y de productos de tienda. Se trata de aislar el efecto de mejorar un atributo físico específico sobre el comportamiento de los consumidores. 
    \begin{itemize}
        \item ¿La cartelería digital aumenta las ventas de la estación?
        \item ¿Qué tipo de promociones funcionan mejor?
        \item ¿Importa el momento del día en que se muestra la promoción?
        \item ¿Influye la ubicación del cartel dentro de la estación?
        \item ¿Genera más cross-selling hacia productos de tienda?
        \item ¿Qué retorno económico tiene la inversión en cartelería digital?
        \item ¿Aumenta el flujo de clientes o el volumen por ticket gracias a un entorno más atractivo y seguro?  
        \item ¿Ciertas mejoras físicas impactan más en las ventas nocturnas o en la captación de cierto perfil de cliente?
    \end{itemize}
    
Este modulo se basa en una secuencia de dos preguntas que se encuentran relacionadas. La carteleria digital esta pensada para mejorar la experiencia del usuario y para aumentar la identificacion entre la marca y el consumidor. Adicionalmente, conidcional en que lo anterior se cumpla, existe una pregunta sobre si estas caracteristicas se traducen en un aumento de ventas. La dificultad en  abordar estas preguntas reside en que se deben separar los efectos por un lado, el impacto en el usuario, y por otro lado el impacto en las ventas. Es decir, es posible que un usuario logre una identificacion mayor con la marca, pero que esto no se traduza en un aumento de combustible u otros serivicios vendidos. 
Una estratgia en donde por ejemplo se encueste a usuarrios antes y despues de la implementacion de la carteleria digital, puede ayudar a medir el impacto en la experiencia del usuario. Sin embargo, uno no podria medir el impacto que esa mejora ha tenido en las ventas. Por otro lado, una estrategia de medicion que intente explicar el aumento de ventas a partir de un aumento en la implementacion,  no necesariamente puede medir el canal por el cual se ha dado ese aumento. Sin conocer el canal, no se puede determinar si el aumento en ventas es por un aumento en la experiencia del usuario, y ahora el consumdidor se fideliza m'as con YPF, o porque la es simplemente un aumento geniuno de la demanda en donde la carteleria digital ha sido implementada, porque  desde antes de la implementacion, esa zona ha tenido un aumento de demanda. Una correcta estrategia de estimacion deberia contemplar las dos cuestiones en simultaneo. 

\subsection{Diseño econométrico}
    Para una estimacion correcta y un analisis integral sobre la implementacion de carteleria digital, se necesita una discusion conceptual con el equipo que se encuentre implementando los conceptos. 
    Los objetivos de la comunicacion comercial dinamica, en general se basan en: 
    \begin{enumerate}
        \item Identificacion del consumidor con la marca.
        \item Experiencia del usuario. 
    \end{enumerate}
    Dentro de los dos puntos anteriores, existen objetivos puntuales. Es decir, en epoca del mundial, una publicidad especifica sobre la seleccion argentina. Una estrategia usual es analizar encuestas a consumidores sobre la relacion con la marca, ejemplo encuestas de calidad de servicio en la estacion de servicvio misma o a traves de \emph{APP YPF}. Dado el nivel desagregado de informacion de la APP, se pueden estimar metricas sobre el impacto de una publicidad especifica en el uso de la Aplicacion. Una primera etapa constaria de individualizar si al cambiar la publicidad, en promedio, la aplicacion se abre mas, o se usa mas tiempo. El hecho de contar con datos a nivel usuario que se pueden linkear con la estacion de servicio donde los clientes cargan combustible genera que se pueda medir que tanto los consumidores interactuan con la aplicacion. 
    La estrategia no se limita a estimar simplemente el impacto de las distintas publicidades en la interaccion del usuario con la aplicacion sino que utiliza ese insumo en una segunda etapa para determinar como la mayor o menor exposicion a la interaccion con la aplicacion genera efectos en la demanda ya sea de combustibles y productos dentro de la tienda. 
    La estrimacion vendria a ser un modelo por variables instrumentales donde la variable instrumental seria la exposicion a la publicidad, que se mide a traves de la interaccion con la aplicacion. La variable endogena seria el aumento en ventas, y el resultado de interes es el impacto de la interaccion con la aplicacion sobre las ventas.
    Los resultados sirven en dos dimensiones. En una primera etapa, se puede rankear el tipo de publicidad a la que los consumidores son mas reactivos, midiendo con una metrica precisa, dada por la interaccion con la aplicacion. POr otro lado, este aumento de interaccion con la aplicacion, puede considerarse exogeno, dado que la publicidad se muestra de manera aleatoria a los consumidores, y entonces estimar de manera causal como la interaccion con la aplicacion se refleja en ventas, ya sea en el momento, o en futuros momentos. 


  \subsection{Datos requeridos}
    
    Para una completa estimacion, se necesitan datos a nivel usuario y anivel estacion de servicio, con frecuencia hora/día. Se necesita informacion sobre los tipos de publicidades mostradas en la carteleria digital.  Por otro lado, se necesitan datos sobre las ventas de combustibles y productos de tienda, con una frecuencia diaria o incluso horaria, para poder medir el impacto de la publicidad en las ventas. Adicionalmente, se necesitan datos sobre la interaccion de los usuarios con la aplicacion, para medir el impacto de la publicidad en la interaccion con la aplicacion. 
    La base de datos ideal tendria una estructura a nivel de estacion de servicio/hora con la cantidad de usuarios que compraron combustible, la cantidad de usuarios que interactuaron con la aplicacion, el tipo de publicidad mostrada en la carteleria digital, y las ventas de combustibles y productos de tienda. 


